\scanpage{120}
\mspage{107}
$\varphi\in TC(\Delta^1)$) alors la flèche $\deleted{\unsure{f_{1,0}}}\widetilde{\varphi}$ dans
le diagramme (6.2.3) est dans $C$ (resp. $TC$).

Plus généralement et plus précisément :

\uline{Lemme} 6.2.5. Soit $M$ une catégorie \inserted{\unsure{stable} par sommes \unsure{amalg.|ind.}},
et $\Phi$ une partie $Q$-close à gauche
(on prendra $\Phi=C$, et $\Phi=TC$, dans le cas
du lemme 6.2.4). \deleted{\unsure{Munissons}} Soit
$\Psi=\Phi_*$, et considérons, pour \unsure{tte} \deleted{\unsure{catég}}
$I$ dans $\mathrm{Cat}$, $\Phi(I)$, $\Psi\langle I\rangle$ dans $M(I)$
\unsure{définis|définies} par
\[
 \begin{aligned}
 &u\colon X\longrightarrow Y\ \inserted{\text{dans }M(I)}\ \text{est dans }\Psi\langle I\rangle\ \text{sss}\\
 &\text{\unsure{ssi|sss}}\quad \forall i\in\operatorname{Ob} I,\quad u(i)\colon X(i)\unsure{\text{---}|\longrightarrow}Y(i)\\
 &\text{est dans }I,
 \end{aligned}\tag*{(6.2.5.1)}
\]
\[
 \Phi(I)=\Psi\langle I\rangle^*.\tag*{(6.2.5.2)}
\]
Prenons le cas $I=\Delta^1\inserted{=\{0\xrightarrow{\unsure{d|\delta}}1\}}$, et soit
$\deleted{\unsure{f}}u\colon X\longrightarrow Y$ une flèche dans $M(\Delta^1)$
correspondant à un carré comm.

\diagram{\begin{tikzcd}[ampersand replacement=\&]
\& X(1)\arrow[r,"{X(d)}"]\arrow[d,"{u(1)}"'] \& X(0)\arrow[d,"{u(0)}"]\\
(6.2.5.3)\& Y(1)\arrow[r,dash,"{Y(d)}"] \& Y(0).
\end{tikzcd}}

Considérons le \unsure{diagramme|codiagramme} correspondant
% archive sheet number: 119

